\documentclass[11pt]{beamer}
\usetheme{Boadilla}
\usepackage[utf8]{inputenc}
\usepackage[danish]{babel}
\usepackage{amsmath}
\usepackage{amsfonts}
\usepackage{amssymb}
\usepackage{graphicx}
\usepackage{pgfplots}
\usepackage{tikz}

\definecolor{UBCblue}{rgb}{0.95, 0.0, 0.1}

%\useoutertheme{miniframes} % Alternatively: miniframes, infolines, split
\useinnertheme{circles}
\usecolortheme[named=UBCblue]{structure}

\author{}
\title{Julesange 2019}
%\setbeamercovered{transparent} 
%\setbeamertemplate{navigation symbols}{} 
%\logo{} 
%\institute{} 
\date{24. december 2019} 
%\subject{} 

\makeatother
\setbeamertemplate{footline}
{
  \leavevmode%
  \hbox{%
  \begin{beamercolorbox}[wd=.4\paperwidth,ht=2.25ex,dp=1ex,center]{author in head/foot}%
    \usebeamerfont{author in head/foot}\insertshorttitle
  \end{beamercolorbox}%
  \begin{beamercolorbox}[wd=.45\paperwidth,ht=2.25ex,dp=1ex,center]{title in head/foot}%
    \usebeamerfont{title in head/foot}\insertsection
  \end{beamercolorbox}
  \begin{beamercolorbox}[wd=.15\paperwidth,ht=2.25ex,dp=1ex,center]{title in head/foot}%
    \usebeamerfont{title in head/foot}
    \insertframenumber{} / \inserttotalframenumber\hspace*{1ex}
  \end{beamercolorbox}}
  \vskip0pt%
}
\makeatletter
\setbeamertemplate{navigation symbols}{}

\makeatletter
\ifbeamer@notesnormals%
    \beamer@paperwidth 16.00cm%
    \beamer@paperheight 9.00cm%
    \geometry{papersize={\beamer@paperwidth,\beamer@paperheight}}
\fi%
\makeatother


\AtBeginSection[]{
  \begin{frame}
  \vfill
  \centering
  %\begin{beamercolorbox}[sep=8pt,center,shadow=true,rounded=true]{title}
  \begin{Huge}
    \insertsectionhead\par%
  \end{Huge}
  %\end{beamercolorbox}
  \vfill
  \end{frame}
}


\begin{document}

\begin{frame}
\titlepage
\end{frame}

%\begin{frame}
%\tableofcontents
%\end{frame}




\newcommand{\xmastree}{\begin{tikzpicture}
\begin{axis}[
    hide axis, clip=false,
    y domain=0:2*pi,
    samples=30,axis equal, view={65}{20}]
\addplot3 [domain=0:4,surf,shader=flat,z buffer=sort,fill=green!50!brown, draw=green!50!black, line join=bevel] 
    ({(1.3*x-floor(x))*x^0.5*cos(deg(y))},
     {(1.3*x-floor(x))*x^0.5*sin(deg(y))},
     {-5*x});
\addplot3 [domain=0:4.5, samples=10, samples y=0, fill=yellow!85!red, draw=yellow!55!red] ({0},{sin(x/5*360) * (1 - x + floor(x)) },{cos(x/5*360) * (1 - x + floor(x) )  + 0.5}) -- cycle;
\end{axis}
\end{tikzpicture}}

\newcommand{\toc}{\begin{frame}{Sangoversigt}
\begin{columns}
	\begin{column}{0.05\textwidth}
	   %test
	\end{column}
	\begin{large}
		\begin{column}{0.65\textwidth}
    	    \tableofcontents
    	\end{column}

	\end{large}
    
    \begin{column}{0.30\textwidth}
    	\begin{center}
    		%\xmastree
    	\end{center}
    \end{column}
\end{columns}



\end{frame}}

%\begin{frame}
%\xmastree
%\end{frame}

\toc

\section{Lille Messias}
\huge
\center
\begin{frame}
En stjerne på himlen

Og flere og flere

I nat føder himlen

Med jordiske veer

I nat lover himlen

At den vil befri os

Maria på æslet

Med lille Messias
\end{frame}

\begin{frame}
Maria skal føde

Sin glæde med smerte

Der banker et hjerte

Bag vinterens hjerte

Selv kulden og mørket

Har håb de vil gi os

Maria på æslet

Med lille Messias
\end{frame}

\begin{frame}
I nat ruger himlen

På glæden og freden

Så varsomt i mørket

Som fuglen i reden

Og freden og glæden

Nu fødes de i os

Maria på æslet

Med lille Messias
\end{frame}

\begin{frame}
I fødslen er glæden

Forbundet med smerten

Dog aldrig var glæden

Så hjemløs i verden

Men stalden og krybben

Er tegn som skal si' os

Maria fandt plads

Da hun fødte Messias
\end{frame}

\toc

\section{En rose så jeg skyde}
\begin{frame}
En rose så jeg skyde

op af den frosne jord

alt som os fordum spåde

profetens trøsteord.

Den rose spired frem

midt i den kolde vinter

om nat ved Betlehem.
\end{frame}

\begin{frame}
Den rose fin og lille

har dejlig duft og skær;

den lyse for os ville

og sprede mørket her;

i sandhed mand og Gud

af syndens nød og pine

han nådig hjalp os ud.
\end{frame}

\begin{frame}
For rosen nu jeg kvæder

omkap med himlens hær:

en jomfru var hans moder,

Maria ren og skær.

I ham brød lyset frem

midt i den mørke vinter

om nat ved Betlehem.
\end{frame}
\toc

\section{Dejlig er jorden}
\begin{frame}
Dejlig er jorden,

prægtig er Guds himmel,

skøn er sjælenes pilgrimsgang!

Gennem de fagre

riger på jorden

gå vi til paradis med sang!
\end{frame}

\begin{frame}
Tider skal komme,

tider skal henrulle,

slægt skal følge slægters gang.

Aldrig forstummer

tonen fra himlen

i sjælens glade pilgrimssang.
\end{frame}

\begin{frame}
Englene sang den

først for markens hyrder,

skønt fra sjæl til sjæl det lød:

Fred over jorden!

Menneske, fryd dig,

os er en evig frelser fød! 
\end{frame}
\toc

\section{På loftet sidder nissen}
\begin{frame}

På loftet sidder nissen med sin julegrød,

sin julegrød, så god og sød,

han nikker, og han spiser, og han er så glad

for julegrød er hans bedste mad.

Men rundt omkring står alle de små rotter,

og de skotter, og de skotter:

''Vi vil så gerne ha' lidt julegodter,''

og de danser, danser rundt i ring.
\end{frame}

\begin{frame}
Men nissefar han truer med sin store ske:

''Vil I nu se at kom' af sted,
for jeg vil ha' min julegrød i ro og fred,
og ingen, ingen vil jeg dele med.''

Men rotterne de hopper, og de danser,

og de svinser, og de svanser,

de kikker efter grøden, og de standser,

og de står om nissen tæt i ring.
\end{frame}

\begin{frame}
Men nissefar han er en lille hidsigprop,

og med sin krop han gør et hop:

''Jeg henter katten, hvis I ikke holder op,
når katten kommer, skal det nok bli' stop.''

Så bliver alle rotterne så bange,

åh, så bange, åh, så bange,

de vender sig og danser et par gange,

og en, to, tre, så er de væk! 
\end{frame}

\toc

\section{Der bor en lille nissemand}
\begin{frame}
Der bor en lille nissemand

på loftets mørke rum, bummelum.

Der bor en lille nissemand

på loftets mørke rum.

Han stamper rundt, han tramper rundt,

han spiser grød, for det er sundt.

Der bor en lille nissemand

på loftets mørke rum.
\end{frame}

\toc

\section{Fire strømper uden fod}
\begin{frame}
Fire strømper uden fod,

en tomat så rød som blod

og en gammel gulerod

ka’ du få for en krone
\end{frame}

\begin{frame}
Sørøver Olsens sorte kat,

og en lille pakke vat,

og den hemmelige skat

ka’ du få for en krone.
\end{frame}

\begin{frame}
Et par vanter, som er blå,

og en bøf med peber på

og halvanden dameskrå

ka’ du få for en krone.
\end{frame}

\begin{frame}
Knæk og bræk med skrin og spæk,

syv kartofler i en sæk

og et blækhus uden blæk

ka’ du få for en krone.
\end{frame}

\begin{frame}
Kalles lysegrønne flag

og et drops med hindbærsmag

og en revnet buksebag

ka’ du få for en krone.
\end{frame}

\toc

\section{Den himmelblå}
\begin{frame}
Se, der' den nye julemand

Mon han har klejner med?

Han vinker til os fra sit rensdyrspand

Mens han flyver afsted
\end{frame}

\begin{frame}
Jesus bli'r født på anden sal

I Magasin du Nord

Han griner, når du trykker på en fodpedal

Si'r goddag til sin mor
\end{frame}

% Begin chorus
\begin{frame}
Nu' det jul, og træet tændes påny
Nu' det jul i land og by
Og guds engle små klæ'r sig festligt på
I den himmelblå
\end{frame}
\begin{frame}
Nu' det jul, og træet tændes påny
Nu' det jul i land og by
Og guds engle små klæ'r sig festligt på
Og synger i sky
\end{frame}
% End chorus

\begin{frame}
Familien er samlet fra nær og fjern
Synet er pyntet med vat
Og udenfor lyser den ledestjerne
I en smuk julenat
\end{frame}

% Begin chorus
\begin{frame}
Nu' det jul, og træet tændes påny
Nu' det jul i land og by
Og guds engle små klæ'r sig festligt på
I den himmelblå
\end{frame}
\begin{frame}
Nu' det jul, og træet tændes påny
Nu' det jul i land og by
Og guds engle små klæ'r sig festligt på
Og synger i sky
\end{frame}
% End chorus

\toc
\section{O Come, All Ye Faithfull}
\begin{frame}
O come, all ye faithful

Joyful and triumphant

O come ye, o come ye to Bethlehem

Come and behold Him

Born the King of Angels!
\end{frame}

\begin{frame}
O come, let us adore Him

O come, let us adore Him

O come, let us adore Him

Christ the Lord
\end{frame}

\begin{frame}
Oh, sing, choirs of angels

Sing in exultation

Oh, come, oh come ye to Bethlehem

Come and behold Him

Born the King of Angels
\end{frame}

\begin{frame}
O come, let us adore Him

O come, let us adore Him

O come, let us adore Him

Christ the Lord
\end{frame}

\toc

\section{Mary's Boy Child}
\begin{frame}
Long time ago in Bethlehem,

so the Holy Bible said

Mary's boy child Jesus Christ,

was born on Christmas Day
\end{frame}

\begin{frame}
Hark, now hear the angels sing,

a king was born today

And man will live for evermore,

because of Christmas Day

Mary's boy child Jesus Christ,

was born on Christmas Day
\end{frame}

\begin{frame}
While shepherds watch their flocks by night

They see a bright new shining star

They hear a choir sing a song,

the music seemed to come from afar
\end{frame}

\begin{frame}
Hark, now hear the angels sing,

a king was born today

And man will live for evermore,

because of Christmas Day

Mary's boy child Jesus Christ,

was born on Christmas Day
\end{frame}

\begin{frame}
When Josef and his wife Mary,

came to Bethlehem that night.

They found no place to birth a child,

not a single room was in sight.
\end{frame}

\begin{frame}
Hark, now hear the angels sing,

a king was born today

And man will live for evermore,

because of Christmas Day

Mary's boy child Jesus Christ,

was born on Christmas Day
\end{frame}

\begin{frame}
At last they found a little nook,

in a stable all forlorn.

And in a manger cold and dark,

Marys little boy was born.
\end{frame}

\begin{frame}
Hark, now hear the angels sing,

a king was born today

And man will live for evermore,

because of Christmas Day

Mary's boy child Jesus Christ,

was born on Christmas Day
\end{frame}

\toc

\end{document}
